We will start with a definition which is not of an $\infty$-category, but a weaker notion.
In particular these will be a simplicial set with composition, but the space of composers isn't necessarily contractible.

\begin{definition}
    A \textbf{composer} is a simplicial set with the extension property against spine inclusions $I^n \to \Delta^n$.
\end{definition}

\begin{example}
    The singular set of a topological space is a composer.
    Given paths $p_i : x_i \to x_{i+1}$, the data of a spine map, we define their composition as the composition of the paths.
    If you aren't convinced this gives an actual $n$-cell, there is some way to see this either purely topologically, or by inductively filling $2$-horns, $3$-horns, etc. until you get a full filler of $\Delta^n$, using 
\end{example}

\begin{definition}
    Two $1$-simplices $f, g : x \to y$ are \textbf{equivalent} if there is a two simplex
    $\sigma : \Delta^2 \to X$ such that
    \[ \sigma_{0,1} = f\]
    \[ \sigma_{0, 2} = g\]
    \[ \sigma_{1,2} = \id_y\]
\end{definition}

\begin{definition}
    Let $X$ be a simplicial set. 
    Define the \textbf{homotopy category} $hX$ to be the category with:
    \begin{enumerate}
        \item Objects: $X_0$
        \item Morphisms: freely generated by $X_1$. Write the free composition by $f \ast g$.
    \end{enumerate}
    quotiented by the relations
    \begin{enumerate}
        \item $s_0(x) = \id_x$
        \item For every $2$-simplex $\Sigma$ with boundary $(f,g,h)$ we have $h = g \ast f$.
        \item If $f \sim f'$ then $f \ast g \sim f' \ast g$ and $g' \ast f \sim g' \ast f'$. 
    \end{enumerate}
\end{definition}

This construction gives a functor $h : \sSet \to \mathrm{Cat}$.

In a composer, the free generation collapses - every formal composite will be equivalent to a single morphism.

\begin{lemma}
    Suppose $X$ is a composer which lifts wrt inner $3$-horn inclusions.
    Let $f$ and $g$ be comoposable morphisms in $X$.
    Then 
    \begin{enumerate}
        \item There exists a composite of $f$ and $g$
        \item The relation of equivalence is an equivalence relation
        \item Any two composites of $f$ and $g$ are equivalent
        \item given a $2$-simplex $\sigma$ with $\sigma = (\id_x, h, h')$ then $h ~ h'$
    \end{enumerate}
    \begin{proof}
        This is an exercise in horn filling!
    \end{proof}
\end{lemma}

Furthermore in such a composer we have a nicer description of its homotopy category:
It is given by the category $\pi(X)$  with objects $X_0$ and morphisms equivalence classes of $X_1$.
Associativity of composition also follows from horn filling.

\begin{definition}
    A simplicial set $X$ is an $\infty$-category if it has fillers for all inner horn inclusions.

    A functor between $\infty$-categories is a map of simplicial sets.
\end{definition}

\begin{definition}
    A morphism in an $\infty$-category is an \textbf{equivalence} if its image in the homotopy category is an isomorphism.

    Equivalently there are two cells $\sigma$ and $\tau$ witnessing whose outer cells witness that $f$ has a right and left inverse.
\end{definition}

\begin{definition}
    Given an $\infty$-category we construct its maximal subgroupoid as a pullback of simplicial sets
    \[\begin{tikzcd}
        \cat{C}^\simeq \ar[r] \ar[d] & \cat{C} \ar[d] \\
        N(h\cat{C}^\simeq) \ar[r] & N(h\cat{C})
    \end{tikzcd}\] 
\end{definition}

\subsection{Examples of infinity categories}

Apart from the usual examples, to cook up more examples we will use simplicially enriched categories.
\begin{definition}
    A  \textbf{simplicial category} is a category enriched in simplicial sets, with the cartesian monoidal structure.

    Explicitly, we have a set $\cat{C}$ of objects, and for each $x,y : \cat{C}$ a simplicial set of morphisms $\cat{C}(x,y)$, with identities and composition.

    By abstract nonsense the category of small simplicial categories is bicomplete.
    Additionally noting that the maps $c : \Set \to \sSet$ and $\pi_0, \mathrm{ev}_0 : \sSet \to \Set$ are monoidal, we obtain maps
    \[ c : \mathrm{Cat} \to \mathrm{Cat}_\Delta \]
    \[ \pi, u : \mathrm{Cat}_\Delta \to \mathrm{Cat} \]
    These form an adjoint triple $\pi \dashv c \dashv u$
\end{definition}

\begin{definition}
    A simplicial functor $f : \cat{C} \to \cat{D}$ of simplicial categories is a \textbf{weak equivalence} if it is a weak equivalence on hom simplicial sets,
    and $\pi(f) : \pi(\cat{C}) \to \pi(\cat{D})$ is essentially surjective.

    This is the ``weak'' version of being fully faithful and essentially surjective.
\end{definition}

\begin{definition}
    A morphism $f$ in a simplicial category $\cC$ is a zero cell of $\cC(x, y)$.
    It is an \textbf{equivalence} if $\pi(f)$ is an isomorphism.
\end{definition}

Now we would like to perform he nerve construction on simplicial categories.
For this we ened to define the simplicially enriched version of $[n]$.

\begin{definition}
    Let $J$ be a finite non-empty linearly ordered set and $i, j : J$.
    Let $P_{i,j} := \{ I \subseteq J \mid i, j \in I \text{ and } k \in I \Rightarrow i \leq k \leq j \}$
    This is partially ordered by inclusion.

    If $i \leq j \leq k$ there is a canonical map
    \[ P_{i, j} \times P_{j,k} \to P_{i,k} \]
    given by the union.

    We can use this to define a simplicial category $\cC[\Delta^J]$ with
    \begin{enumerate}
        \item Objects $J$
        \item morphisms $i \to j$ given by the simplicial set $N(P_{i,j})$.
    \end{enumerate}
    Composition is induced by the canonical map.
\end{definition}

We can explicitly describe these homsets, since for all $n$ we have $P_{0,n} \simeq [1]^{n-1}$, 
by taking the indicator of the subset.
We can also see $P_{i, j} \simeq P_{0, j-i}$. 